% ============================================================
% CAPITOLUL 1 — INTRODUCERE
% ============================================================

\chapter{Introducere}
\label{ch:introducere}

\section{Context și motivație}
\label{sec:context}

Bolile infecțioase rămân una dintre formele prin care vulnerabilitatea unui sistem
sanitar devine vizibilă la scară socială. Pandemia de gripă din 1918, epidemia de
SARS din 2003 și pandemia de COVID-19 declanșată în 2019 arată că propagarea unui
agent patogen nu poate fi înțeleasă doar ca fenomen clinic, ci și ca proces
populațional, dependent de contacte, mobilitate, raportare și capacitatea de
răspuns instituțional \cite{andersonmay1991, hethcote2000}.

Modelele matematice compartimentale oferă un limbaj riguros pentru descrierea
acestor procese. Ele reduc populația la stări epidemiologice relevante și descriu
tranzițiile dintre stări prin ecuații diferențiale. Modelul SIR
(Susceptibili--Infectați--Recuperați), introdus de Kermack și McKendrick în 1927
\cite{kermack1927}, rămâne un reper important deoarece permite o analiză
matematică transparentă. Totuși, pentru epidemii în care testarea, diagnosticul și
severitatea clinică modifică puternic raportarea, un model cu trei compartimente
este prea agregat. Din acest motiv, lucrarea folosește SIDARTHE ca model clinic
principal, păstrând modelele SIR, SIS, SEIR și SIRS ca repere secundare pentru
comparații și verificări.

În practica actuală, problema nu se oprește la scrierea ecuațiilor. Datele provin
din surse eterogene, sunt actualizate frecvent și conțin corecții, goluri sau
inconsistențe administrative. În același timp, implementarea numerică poate fi
realizată în limbaje diferite, fiecare cu avantaje și costuri proprii. Alegerea
dintre R, Julia și Python nu ține doar de sintaxă, ci influențează timpul de execuție,
consumul de memorie, modul de integrare cu infrastructura și reproductibilitatea
analizei.

Pornind de la aceste observații, lucrarea propune un sistem e-medicine în care
modelul SIDARTHE este integrat într-un flux complet: ingestie de date
epidemiologice, execuție asincronă a simulărilor, stocare a rezultatelor și
vizualizare în interfață web. Modelele SIR, SIS, SEIR și SIRS rămân utile ca
repere mai simple, prin care se poate observa ce se pierde sau ce se câștigă
atunci când modelul introduce compartimente suplimentare.

\section{Obiectivele lucrării}
\label{sec:obiective}

Obiectivul principal este proiectarea unui sistem care leagă explicit modelarea
matematică de implementarea software. Modelul SIDARTHE este tratat ca model de
lucru principal, deoarece separă cazurile nedepistate, diagnosticate,
simptomatice, severe, vindecate și decedate. Familia
SIR/SIS/SEIR/SIRS este analizată în continuare, dar rolul său devine acela de
reper teoretic și computațional, util pentru înțelegerea proprietăților de bază
ale modelelor compartimentale.

Al doilea obiectiv este implementarea aceluiași contract de simulare în R, Julia
și Python, folosind bibliotecile uzuale ale fiecărui ecosistem pentru rezolvarea
ecuațiilor diferențiale. Comparația urmărește atât compatibilitatea rezultatelor,
cât și costurile practice ale execuției: timp wall-clock, timp CPU, memorie și
integrare cu baza de date.

Al treilea obiectiv este definirea unei metodologii explicite de comparație și
calibrare. Modelele nu sunt evaluate doar vizual, ci prin operatori de observație,
metrici de eroare precum MAE și RMSE, separare între fereastra de calibrare și
fereastra de test și raportarea costului software asociat fiecărei rulări.

Un obiectiv suplimentar este integrarea unei extensii predictive bazate pe LSTM.
Aceasta este tratată ca model de serie temporală pentru predicții numerice pe
termen scurt, implementat prin runner dedicat, și ca posibilă corecție a
reziduurilor produse de modelele compartimentale calibrate.

În paralel, lucrarea construiește infrastructura necesară rulării acestor
experimente într-un cadru accesibil prin web. API-ul Go, baza de date
YugabyteDB, coada Redis și cele două interfețe Vue.js sunt tratate ca părți ale
aceluiași instrument de cercetare, nu ca module izolate.

\section{Contribuții originale}
\label{sec:contributii}

Contribuția lucrării constă în combinarea unei analize matematice clasice cu un
sistem software funcțional, capabil să ruleze modele compartimentale cu structuri
de stare diferite în aceleași medii de execuție. Partea matematică fixează
ipotezele și limitele modelelor, partea de implementare arată cum se schimbă
costul numeric între limbaje, iar partea de arhitectură oferă un cadru în care
datele, simulările și benchmark-urile pot fi gestionate unitar.

O a doua contribuție este metodologia de colectare a metricilor. Fiecare runner
marchează explicit stările unei simulări, scrie rezultatele în aceeași schemă de
bază de date și salvează metrici comparabile. Această disciplină de integrare este
necesară pentru a evita comparațiile informale între limbaje, care pot fi ușor
influențate de diferențe de I/O, inițializare sau serializare.

O altă contribuție este separarea între comparația epidemiologică și comparația
software. Lucrarea definește operatori de observație pentru a transforma stările
interne ale modelelor în serii comparabile cu datele reale și introduce un cadru
de calibrare care poate fi extins ulterior într-un job dedicat în aplicație.

În fine, sistemul web rezultat permite folosirea modelelor într-un flux apropiat
de un scenariu real de analiză colaborativă: crearea unui studiu, acordarea
rolurilor, încărcarea unui set de date, alegerea modelului și a parametrilor,
execuția asincronă, vizualizarea traiectoriilor și compararea rulărilor.

\section{Structura lucrării}
\label{sec:structura}

Lucrarea este organizată în zece capitole.

\textbf{Capitolul~\ref{ch:epidemiologie}} (\emph{Epidemiologie Computațională și
Modelarea Matematică}) prezintă fundamentele epidemiologiei matematice. Sunt
introduse modelele compartimentale (SIR, SIS, SEIR, SIRS), apoi este prezentat
modelul SIDARTHE ca model clinic principal. Modelul SIR este păstrat pentru
analiza matematică de bază, deoarece permite discutarea clară a pozitivității,
conservării populației și pragului epidemic $\Rzero$. Capitolul include și o
trecere în revistă a metodelor numerice utilizate pentru rezolvarea sistemelor de
ecuații diferențiale ordinare.

\textbf{Capitolul~\ref{ch:tehnologii}} (\emph{Tehnologii și Instrumente Utilizate})
explică de ce au fost alese limbajele R, Julia și Python, baza de date
YugabyteDB, orchestrarea containerelor prin Kubernetes, backend-ul Go și
frontend-urile Vue.js.

\textbf{Capitolul~\ref{ch:arhitectura}} (\emph{Arhitectura Sistemului E-Medicine})
descrie arhitectura globală a sistemului, cerințele funcționale și non-funcționale,
schema bazei de date distribuite, arhitectura API-ului backend și structura
componentelor frontend.

\textbf{Capitolul~\ref{ch:implementare}} (\emph{Implementarea modelelor
compartimentale}) detaliază implementările practice ale modelelor în fiecare limbaj de
programare, inclusiv seturile de date utilizate, parametrii, schema generică de
rezultate și interfața de integrare cu API-ul Go.

\textbf{Capitolul~\ref{ch:lstm}} (\emph{Predicție Epidemiologică folosind Rețele
LSTM}) prezintă rolul unei componente LSTM pentru prognoza seriilor temporale,
formularea problemei de predicție, integrarea cu modelele compartimentale și
evaluarea unui model hibrid compartimental--LSTM.

\textbf{Capitolul~\ref{ch:comparatie_calibrare}} (\emph{Comparația Modelelor,
Analiză Numerică și Calibrare}) introduce tabelul comparativ al modelelor,
operatorii de observație, metricile MAE/RMSE și formularea calibrării
parametrilor ca problemă de optimizare.

\textbf{Capitolul~\ref{ch:evaluare}} (\emph{Evaluarea Software, Performanței și
Scalabilității}) prezintă metodologia de benchmarking, matricea experimentală,
metricile colectate, verificările de regresie numerică și evaluarea API-ului și a
bazei de date.

\textbf{Capitolul~\ref{ch:aplicatie}} (\emph{Aplicația Web E-Medicine})
detaliază implementarea aplicației web, documentând backend-ul Go, aplicația
principală Vue.js și consola Vue.js pentru super-admini, cu descrierea
funcționalităților și a fluxurilor de utilizare.

\textbf{Capitolul~\ref{ch:concluzii}} (\emph{Concluzii și Direcții Viitoare})
sintetizează contribuțiile lucrării, discută limitele sistemului și identifică
direcții de dezvoltare ulterioară.

Lucrarea se încheie cu \textbf{Bibliografie} și \textbf{Anexe} care conțin
codul sursă complet al implementărilor, schema bazei de date, configurațiile
Kubernetes și documentația API-ului REST.
